%%%%%%%%%%%%%%%%%%%%%%%%%%%%%%%%%%%%%%%%%%%%%%%%%%%%%%%%%%%%%%%%%%%%%%%%%%
%%%%%%%%%%%%%%%%%%%%%%%%%%%% MP10 %%%%%%%%%%%%%%%%%%%%%%%%%%%%%%%%%%%%%%%%
%%%%%%%%%%%%%%%%%%%%%%%%%%%%%%%%%%%%%%%%%%%%%%%%%%%%%%%%%%%%%%%%%%%%%%%%%%

\section{Variável MP10}

Esta seção reúne os resultados completos obtidos para a variável MP10 durante as etapas experimentais realizadas nesta pesquisa.

\subsection{Resultados das cinco rodadas experimentais}

Esta subseção apresenta os resultados obtidos nas cinco rodadas experimentais realizadas no cenário com falhas simuladas para a variável MP10. Os valores apresentados correspondem diretamente às saídas produzidas automaticamente pelo sistema desenvolvido, preservando integralmente os resultados utilizados para o cálculo das médias, dos desvios padrão e das demais análises estatísticas apresentadas no Capítulo~3.

\clearpage

\begingroup
\footnotesize
\setlength{\tabcolsep}{2pt}
\renewcommand{\arraystretch}{0.95}
\setlength{\LTleft}{0pt}
\setlength{\LTright}{0pt}

\input{Conteudo/Apendices/tabela_anexo_resultados_mp10.tex}

\endgroup

\clearpage

Os resultados apresentados nesta subseção permitem verificar individualmente o desempenho obtido em cada rodada experimental, constituindo a base para os resultados consolidados da variável MP10 discutidos ao longo desta dissertação.

\subsection{Resultados completos de todas as configurações avaliadas}

Durante o processo de otimização dos modelos foram avaliadas diferentes combinações entre algoritmos de aprendizado de máquina, tamanhos de janela e avanços da janela deslizante.

As tabelas apresentadas nesta subseção reúnem todas as configurações executadas para a variável MP10, permitindo a rastreabilidade dos experimentos realizados e possibilitando a reprodução dos resultados apresentados nesta dissertação.

\clearpage

\begingroup
\tiny
\setlength{\tabcolsep}{1.2pt}
\renewcommand{\arraystretch}{0.90}
\setlength{\LTleft}{0pt}
\setlength{\LTright}{0pt}

\input{Conteudo/Apendices/tabela_apendice_todas_configuracoes_mp10.tex}

\endgroup

\clearpage

Os resultados consolidados apresentados no Capítulo~3 foram obtidos diretamente a partir das configurações reunidas nesta subseção, selecionando-se aquelas que produziram o melhor desempenho para cada tipo de falha avaliada.
